\chapter{Introducción}
\label{ch:introduccion}

\paragraph*{}El estudio del comportamiento académico del alumnado universitario ha ido ganando peso en los últimos años dentro del ámbito de la educación superior. No es una preocupación menor. Detrás de cada estudiante que se abandona hay algo más que una baja en el sistema: hay una trayectoria truncada, recursos que no llegan a ningún puerto y, con demasiada frecuencia, una señal de alerta que nadie supo detectar a tiempo. Por eso, entender qué factores influyen en la continuidad académica no es un ejercicio teórico, sino algo con consecuencias muy reales para personas e instituciones. Y es que cuando una universidad pierde a un estudiante a mitad del camino, no solo falla el sistema de retención, sino que se pierde también una oportunidad de intervenir que, con las herramientas adecuadas, podría haberse aprovechado.

\paragraph*{}En ese escenario, el análisis de datos se ha convertido en un aliado cada vez más valioso. No porque los números lo expliquen todo, sino porque permiten ver patrones que a simple vista resultan invisibles como qué perfiles tienen más riesgo, en qué momentos del curso se concentran las bajas o qué variables marcan la diferencia entre quien continúa y quien no. Aplicar estas técnicas al abandono universitario abre una vía verdaderamente útil para pasar de la intuición a la evidencia, y de la evidencia a la acción.

\paragraph*{}Este Trabajo Fin de Grado nace precisamente de esa inquietud. Toma como punto de partida un dashboard interactivo desarrollado en un trabajo previo, orientado al análisis del abandono universitario a partir de datos reales de estudiantes. Aquel sistema fue un primer paso valioso, pero dejó abiertas varias puertas: modelos estadísticos limitados, una arquitectura mejorable y un enfoque todavía demasiado descriptivo. El objetivo de este Trabajo de Fin de Grado es dar el siguiente salto y construir sobre esa base algo más ambicioso.

\paragraph*{}Para ello, se plantean mejoras en varios frentes. A nivel técnico, se refuerza la arquitectura de la aplicación para ganar en escalabilidad, rendimiento y estabilidad. A nivel analítico, se incorporan nuevas técnicas estadísticas y modelos con mayor capacidad predictiva, capaces de identificar patrones de comportamiento y anticipar posibles situaciones de abandono antes de que sea demasiado tarde para intervenir. Y a nivel de usabilidad, se desarrollan herramientas de visualización más claras e interpretables, se añade soporte multilingüe y se amplían las opciones de exportación de resultados, pensando siempre en quienes, al final del día, tienen que tomar decisiones a partir de esa información.

\paragraph*{}El objetivo final además de mostrar información es transformarla en conocimiento útil. Una buena visualización puede resultar atractiva y estar bien construida técnicamente, pero su verdadero valor aparece cuando ayuda a entender qué está ocurriendo y, sobre todo, qué podría hacerse a continuación. Eso es lo que se busca con este proyecto: contribuir a que los responsables académicos dispongan de una herramienta real para comprender, anticipar y, en la medida de lo posible, prevenir el abandono en el entorno universitario.